\Article{12}
Da die Funktion $\varphi r$ für jeden merklichen Wert von $r$ unmerklich ist und mit wachsendem $r$ stetig abnimmt, folgt unmittelbar, dass das Integral $\int rr\varphi r.\mathrm{d}r$, von einem merklichen Wert bis zu einem größeren erstreckt, ebenfalls unmerklich bleibt, sofern der größere Wert in dem Bereich liegt, in dem wir Versuche anstellen können. Aus dieser Eigenschaft allein dürften wir jedoch keineswegs schließen, dass das Integral bei der Ausdehnung der Integration auf ein beliebig großes Intervall unmerklich bleibt. Die Rechnungen des berühmten \textsc{Laplace} sind zwar so formuliert, dass sie eine solche Annahme einschließen. Solange aber die Natur der Funktion $\varphi r$ unbekannt ist, scheint es ratsamer, auf alle hypothetischen Annahmen zu verzichten, die wir entbehren können. Da uns also die Wahl der Integrationskonstanten in $\int rr\varphi r.\mathrm{d}r=-\psi r$ freisteht, genügt es anzunehmen, dass sie so gewählt ist, dass $\psi r=0$ für einen beliebigen merklichen Wert von $r$ innerhalb des Bereichs der Abmessungen von Körpern, an denen wir Versuche anstellen können. Auf diese Weise ist $\psi r$ für jeden anderen derartigen Wert stets unmerklich (positiv für einen kleineren, negativ für einen größeren). Das schließt jedoch keineswegs aus, dass sie für einen unmerklichen Wert von $r$ merklich werden kann. Hinzuzufügen ist allerdings, dass zur Erklärung der Erscheinungen bei unbegrenzt abnehmendem Abstand $r$ der Wert von $\psi r$ stets endlich bleiben muss, dass also $\psi0$ eine endliche Größe sein muss. Übri%
\newcommand{\NearTau}{\mathord{\text{\fontspec{Times New Roman}τ}}}%
\expandafter\def\csname HistoricalNoteText6\endcsname{Wenn man die Erklärung der Lichterscheinungen nach der Emissionstheorie gelten lässt, hängt die Brechung von der molekularen Anziehung ab, die die Teilchen des durchsichtigen Körpers auf die Lichtteilchen ausüben, und das Brechungsverhältnis vom Wert von $\psi0$, sodass gilt
\[
\frac{c\psi0}{g}=\frac{(nn-1)kk}{8\pi^3l}
\]
wobei $l$ die Länge eines Pendels bezeichnet, das für einen Schwung eine Sekunde braucht, $k$ den Weg des Lichts im Vakuum in einer Sekunde und $n$ das Verhältnis des Sinus des Einfallswinkels zum Sinus des Brechungswinkels. Damit ist für Wasser $\frac{c\psi0}{g}$ zweitausendmal so groß wie der mittlere Abstand der Sonne von der Erde.}%
\SourceLeaf{52}gens hat $\frac{c\psi r}{r}$ offenbar dieselbe Dimension wie die Schwerkraft $g$, oder anders gesagt, $\frac{c\psi r}{g}$ ist eine Länge. Damit ist $\frac{c\psi0}{g}$ eine bestimmte Länge (je nach der Natur der Körper, auf deren Anziehungskräfte sich die Funktion $fr$ bezieht). Man darf zwar vermuten, dass sie ungeheuer groß ist, kann sie aber in bestimmten Fällen kaum auch nur näherungsweise angeben, jedenfalls nur unter hypothetischen Annahmen\HistoricalNote{6}.

\Article{13}
Ganz entsprechend nehmen wir bei der Integration $\int\psi r.\mathrm{d}r=-\theta r$ an, dass die Konstante so gewählt ist, dass $\theta r=0$ für einen beliebigen Wert von $r$ innerhalb des Bereichs der Werte, für die wir Versuche anstellen können. Auf diese Weise ist $\theta r$ für jeden derartigen merklichen Wert von $r$ unmerklich, auch wenn sie für einen unmerklichen Wert merklich werden kann. Offenbar gibt $\frac{c\theta r}{g}$ die Fläche einer zweidimensionalen Figur an und somit $\frac{\theta r}{\psi r}$ eine Länge. Nun muss aber $\frac{\theta0}{\psi0}$ eine Länge von unmerklicher Größe sein, wie wir folgendermaßen zeigen. Da $\psi r$ von $r=0$ an stetig abnimmt, und zwar so rasch, dass sie bereits unmerklich geworden ist, sobald $r$ einen merklichen Wert erreicht, muss derjenige Wert von $r$, für den $\psi r=\frac12\psi0$ gilt, unmerklich sein. Bezeichnen wir ihn mit $\rho$. Betrachten wir das Integral $\int(\psi0-\psi r)\mathrm{d}r$, das bei Ausdehnung von $r=0$ bis $r=R$ den Wert $=R\psi0-\theta0+\theta R$ annimmt. Dieses Integral ist offenbar größer als dasselbe von $r=\rho$ bis $r=R$ erstreckte Integral, und dieses wiederum größer als das Integral $\int(\psi0-\psi\rho)\mathrm{d}r$ zwischen denselben Grenzen. Da das letzte Integral den Wert $=(\psi0-\psi\rho)(R-\rho)=\frac12\psi0.(R-\rho)$ annimmt, gilt allgemein für jeden Wert von $R$ (größer als $\rho$)
\[
R\psi0-\theta0+\theta R>\frac12\psi0.(R-\rho)
\]
Setzen wir nun für $R$ den Wert des Bruchs $\frac{\theta0}{\psi0}$ ein, so ergibt diese Beziehung
\[
\theta R>\frac12\psi0.(R-\rho)
\]
was absurd wäre, wenn $R$ eine merkliche Größe wäre.

\SourceLeaf{53}Trotz der ungeheuren Größe von $\psi0$ steht also nichts dem entgegen, dass $\theta0$ eine recht mäßige Größe sein kann, die mit den Abmessungen der in den Versuchen verwendeten Körper vergleichbar ist.

\Article{14}
Es bleibt zu untersuchen, was aus dieser Eigenschaft der Funktion $\theta$ für das Integral (I)
\[
\int\int\frac{\mathrm{d}t.\mathrm{d}T.\cos q.\cos Q.\theta(\mathrm{d}t,\mathrm{d}T)}{(\mathrm{d}t,\mathrm{d}T)^2}
\]
folgt. Diese Untersuchung muss mit einer einfacheren beginnen. Wir betrachten auf einer der beiden Flächen einen bestimmten Punkt $\mu$ und werten das Integral (II)
\[
\int\frac{\mathrm{d}t.\cos q.\cos Q.\theta(\mu,\mathrm{d}t)}{(\mu,\mathrm{d}t)^2}
\]
aus, das über die ganze Fläche $t$ zu erstrecken ist. Dabei bezeichnet $Q$ den Winkel zwischen zwei vom Punkt $\mu$ ausgehenden Geraden, von denen die eine auf das Element $\mathrm{d}t$ gerichtet und die andere fest ist. Dagegen bezeichnet $q$ den Winkel zwischen zwei vom Punkt des Elements $\mathrm{d}t$ ausgehenden Geraden, von denen die eine auf den Punkt $\mu$ gerichtet ist und die andere senkrecht zum Element nach außen weist.

Zunächst bemerken wir, dass, wenn der Punkt $\mu$ in merklichem Abstand von der Fläche $t$ liegt, die Werte aller $\theta(\mu,\mathrm{d}t)$ unmerklich sind. In diesem Fall ist also das ganze Integral (II) unmerklich. Dieses Integral kann daher nur insoweit einen merklichen Wert annehmen, als die Fläche $t$ Teile aufweist, die in unmerklichem Abstand vom Punkt $\mu$ liegen. Offenbar genügt es, das Integral (II) über diese Teile zu erstrecken und alle Teile in merklichen Abständen zu vernachlässigen.

Ferner ersetzen wir $\frac{\mathrm{d}t.\cos q}{(\mu,\mathrm{d}t)^2}$ durch $\pm\mathrm{d}\Pi$. Dabei bezeichnet $\mathrm{d}\Pi$ auf der Kugelfläche mit Radius $=1$ um den Mittelpunkt $\mu$ dasjenige Element, auf das das Element $\mathrm{d}t$ vom Punkt $\mu$ aus gesehen projiziert wird. Es gilt das obere beziehungsweise das untere Vorzeichen, je nachdem, ob das Element $\mathrm{d}t$ dem Punkt $\mu$ seine Außen- oder seine Innenseite zuwendet. Auf diese Weise erhält das Integral (II) die Form
\[
\int\pm\mathrm{d}\Pi.\cos Q.\theta(\mu,\mathrm{d}t)
\]
und es ist ersichtlich, dass der Wert dieses Integrals nur insoweit merklich werden kann, als die Elemente $\mathrm{d}\Pi$, die zu unmerklichen Abständen $(\mu,\mathrm{d}t)$ gehören, auf der Kugelfläche ein Gebiet von merklicher Größe ausfüllen.

Daraus lässt sich leicht schließen, dass unser Integral im Allgemeinen auch dann unmerklich bleibt, wenn der Punkt $\mu$ auf der Fläche $t$ selbst liegt. Denn offenbar liegen die Pro\SourceLeaf{54}jektionen aller Elemente $\mathrm{d}t$ in unmerklichem Abstand vom Punkt $\mu$ ihrerseits in unmerklichem Abstand von dem Großkreis, den die Tangentialebene an die Fläche $t$ im Punkt $\mu$ auf der Kugelfläche bildet. Drei Fälle sind auszunehmen.

1) Die Krümmungsradien der Fläche $t$ im Punkt $\mu$ sind von unmerklicher Größe.

2) Die Stetigkeit der Krümmung ist im Punkt $\mu$ oder in unmerklichem Abstand von ihm unterbrochen (vgl.\ \textit{Disquiss. gen. circa superficies curvas} Art.\ 3).

3) Die Fläche $t$ weist einen anderen Teil in unmerklichem Abstand vom Punkt $\mu$ auf, etwa wenn die Dicke des Raumbereichs $s$ in der Nähe dieses Punktes unmerklich ist. Diesen Fall kann man übrigens zu demjenigen rechnen, den wir im nächsten Artikel behandeln.

